\runninghead{Forbearance and Unveiling}
\properlane{hidden-presence-and-forbearing-patience}{Divine Forbearance and the Unveiling of an Already-Present Lord}

Why does God appear distant or silent while human hearts grow numb and the master of the house seems away on a long journey? That question rises from four passages of the First Sunday of Advent in Year~B. Isaiah complains that God has let Israel wander and has hidden his face (Isa.~63:17; 64:7~Vg.); the psalmist begs God to wake his sleeping strength and shine his face (Ps.~80[79]:3b--4); Paul tells the Corinthians they are waiting for the unveiling of Jesus Christ (1~Cor.~1:7); and Mark records a parable of a householder who leaves his house for a far country without announcing his hour of return (Mark~13:33--37). St.~Jerome on Isaiah, St.~John Chrysostom on 1~Corinthians, and St.~Hilary of Poitiers and St.~Bede on Mark each address a passage where God appears silent or absent; taken together, their arguments illuminate the Advent Mass as a vigil before a Lord who is already present.

\subsection{When Forbearance Feels Like Absence}

At the threshold of the Mass, the assembly prays the Entrance Antiphon \emph{Ad te levavi} (Ps.~25[24]:1--3): the psalmist lifts his soul to God and asks not to be put to shame, resting on the promise that those who endure in waiting (\emph{qui sustinent te}) will not be confounded (Augustine, \work{Enarrationes in Psalmos}, Ps.~24 §§1--3; Bellarmine, Ps.~24, explanations~1--3). With the \work{Gloria} omitted (GIRM~53), the Collect (\emph{Da, quaesumus, omnipotens Deus}; modern text withheld) opens Advent prayer; in its ancient Gelasian antecedent (II.lxxxiv), the assembly asks almighty God for the will to go forth fittingly in righteous deeds to meet Christ at his coming --- a petition for active readiness during the long wait.

Isaiah~63:17 raises the sharpest accusation of the day: ``Why hast thou made us to err, O Lord, from thy ways: why hast thou hardened our heart?'' In \work{Commentarii in Isaiam} XVII, St.~Jerome refuses to make God the cause of sin and traces that feeling of abandonment to God's \emph{nimia patientia} (``excessive forbearance''): because God delays punishment in order to wait for repentance, his slowness looks from below like indifference and lets unchastened hearts grow cold. Yet Jerome points to the abrupt grammatical leap in the next verses: after crying \emph{Utinam dirumperes caelos, et descenderes} (``O that thou wouldst rend the heavens, and wouldst come down,'' Isa.~64:1~Vg.), the speakers break directly into second-person past address --- \emph{Descendisti, Verbum caro factum es, et habitasti in nobis} (``Thou didst come down; thou, the Word, wast made flesh and didst dwell among us,'' PL~24 col.~646) --- as a prayer already answered in the Incarnation, melting hostile powers before his face.

Near the end of the First Reading, Israel laments, ``thou hast hid thy face from us'' (Isa.~64:7~Vg., \emph{Abscondisti faciem tuam}). Commenting on that very verse (PL~24 col.~650), Jerome immediately cites Psalm~80(79):4 --- \emph{ostende faciem tuam, et salvi erimus} (``shew us thy face, and we shall be saved'') --- the exact refrain of the Responsorial Psalm. St.~Augustine (\work{Enarrationes in Psalmos}, Ps.~79 §3) asks why the psalm begs God to illumine his face when God's face cannot grow dark, and replies that Christ veiled it behind a ``cloud of flesh'' and weakness so that many failed to recognize him until the Resurrection and Ascension revealed who he was. Where a U.S.\ parish sings either of the permitted seasonal Advent Common Psalms (Ps.~25[24] or Ps.~85[84], GIRM~61), the congregation remains in the same petitions of steadfast endurance and desired manifestation.

\subsection{The Unveiling of an Already-Present Lord}

St.~John Chrysostom (\work{Homilies on First Corinthians}, Homily~II, §6) draws a closely related inference from a single Greek noun in the Second Reading. Paul tells the Corinthians they lack no spiritual gift while awaiting the \emph{apokalupsis} (Latin \emph{revelatio}, ``unveiling'') of our Lord Jesus Christ (1~Cor.~1:7). Chrysostom argues that Paul's choice of \emph{revelation} implies Christ is present even now though unseen, and will appear openly at the end. Waiting in Advent is therefore directed toward the lifting of a veil over a Lord who already sustains his Church. Before the Gospel, the Gospel Acclamation \emph{Ostende nobis, Domine} (Ps.~85[84]:8) repeats that verb of showing, asking God to show his mercy and give his salvation --- ``Thy Christ,'' in Augustine's concise phrase (\work{Enarr.}, Ps.~84 §6).

Mark~13:33--37 (following the omitted verse~13:32) then raises the obvious counter-question: if Christ is already present, why does he withhold the hour of his return and compare himself to a householder departing for a far country? St.~Hilary of Poitiers (\work{De Trinitate} IX.65, 67) answers Arian readers who took Mark~13:32 and Matthew~24:36 as evidence of creaturely ignorance in the Son. Pairing Christ's silence with Matthew's watchfulness sayings (Matt.~25:13; Matt.~24:44, 46), Hilary argues that Christ withholds a fixed date as a deliberate economy: certain knowledge of the schedule would breed negligence, whereas suspense stretches faith into continuous vigilance --- a reasoning St.~Thomas Aquinas applies directly to Mark~13:33 in the \work{Catena Aurea in Marcum}. St.~Bede (\work{In Marci Evangelium Expositio} IV.13) keeps Hilary's phrase of poised expectation (\emph{pendulae exspectationis incerto}) and resolves the householder's journey by citing St.~Gregory the Great (\work{Homiliae in Evangelia}, Hom.~9): because earth is the native place of human flesh, Christ went into a ``far country'' when he carried his risen humanity to heaven at the Ascension. He left his Church bodily (\emph{corporaliter}), without ever withdrawing the safeguard of his divine presence.

\subsection{Hidden Nearness at the Altar and the Weight of Sin}

At the altar, that unseen presence becomes sacramental nearness. In the Prayer over the Offerings (\emph{Suscipe, quaesumus, Domine}; modern text withheld), the ancient Veronense antecedent (XXVIIII.372) places gifts drawn from God's own benefits before the Lord so that temporal devotion may become the reward of eternal redemption. At the Seasonal Preface of Advent (GIRM~364; modern wording withheld), Eucharistic thanksgiving stands between the two descents St.~Cyril of Jerusalem contrasts in \work{Catechesis}~15 (§§1--2): Christ's first coming ``unobserved, like rain on a fleece,'' and his second open in glory. Between those two visible comings, St.~Bernard of Clairvaux describes a hidden coming into those who love Christ and keep his word (\work{Sermones in Adventu}~3.4, citing John~14:23), a middle advent in spirit and virtue bridging the first coming in flesh and the last in glory (Sermo~5.1, on Ps.~68[67]:14). At Holy Communion, the Communion Antiphon \emph{Dominus dabit benignitatem} (Ps.~85[84]:13) announces that the Lord gives goodness --- interior sweetness, in Augustine's exposition (Ps.~84 §12) --- so that our earth yields fruit even before the final unveiling. The Prayer after Communion (\emph{Prosint nobis}; modern text withheld) closes the Communion Rite, and the Concluding Blessing and dismissal (GIRM~90) send the assembly back into the week ahead.

This interpretation carries a clear exegetical and pastoral trade-off. Centering the Sunday on divine forbearance and hidden presence adopts St.~Jerome's Incarnational reading of \emph{Descendisti} in Isaiah~64:3b~Vg.\ (Heb.~64:2b) rather than St.~Thomas Aquinas's literal reference to God's deliverance from Egypt and at the cliffs of Arnon. If pressed without balance, emphasizing God's hiddenness can also soften the sharp confession of human sin and stained righteousness in Isaiah~64:5--7~Vg.\ (Heb.~64:4--6) and the rebuke of Corinthian factions that follows immediately in 1~Corinthians~1:10.

\subsection{The Four Senses of Scripture}

\begin{fourSenses}
\item[Literal.] Israel cries out from perceived divine silence and begs God to rend the heavens (Isa.~63:16b--64:8~Vg.), while the psalmist pleads for the Shepherd of Israel to shine his face upon his ravaged vineyard (Ps.~80[79]); Paul reminds the Corinthians that they already enjoy God's gifts while awaiting Christ's manifestation (1~Cor.~1:3--9); Jesus on the Mount of Olives warns his disciples not to mistake his coming departure and withheld timetable for permission to sleep (Mark~13:33--37).
\item[Allegorical.] God's sudden descent (\emph{Descendisti}, Jerome on Isa.~64:3b~Vg.\ [Heb.~64:2b]), the lifting of the veil of flesh (Augustine on Ps.~80[79]:4), and the householder traveling to a far country (Bede on Mark~13:34, drawing on Gregory's Homily~9 on Matt.~25:14--30) signify the Incarnation and Ascension, in which Christ conceals his majesty beneath humanity and sacrament while remaining divinely present in his Church (Chrysostom on 1~Cor.~1:7).
\item[Moral.] Because divine forbearance can look like absence (Jerome on Isa.~63:17) and a known schedule breeds laxity (Hilary, \work{De Trinitate} IX.67, on the withheld day of Christ's coming), believers resist spiritual numbness through watchful prayer, refusing to mistake God's patience for indifference.
\item[Anagogical.] Every hidden coming in flesh, word, and Eucharist moves toward the \emph{apokalupsis} of 1~Corinthians~1:7 --- the open unveiling of the Lord's face (Ps.~80[79]:4) when Christ appears in glory with the angels (Cyril of Jerusalem, \work{Catechesis}~15 §1).
\end{fourSenses}
